\documentclass[a4paper,12pt]{article}
\usepackage[utf8]{inputenc}
\usepackage[T1]{fontenc}
\usepackage[ngerman]{babel}
\usepackage{geometry}
\usepackage{amsmath}
\usepackage{amssymb}
\usepackage{booktabs}
\geometry{margin=2.5cm}

\begin{document}

% ================================================================
\section*{Choose a Transformer: Fourier or Galerkin}

\textbf{Autor:} Shuhao Cao \\
\textbf{Institution:} Department of Mathematics and Statistics, Washington University in St.\ Louis \\
\textbf{Konferenz:} NeurIPS 2021 \\
\textbf{arXiv:} 2105.14995

% ================================================================
\subsection*{1\quad Problemstellung und Motivation}

Das Paper überträgt den Selbst-Aufmerksamkeitsmechanismus aus dem Transformer-Modell
(\emph{Attention Is All You Need}) erstmals auf das datengetriebene Operatorlernen für
partielle Differentialgleichungen. Ausgangspunkt ist die Beobachtung, dass ein trainierter
Operatorlerner -- im Gegensatz zu klassischen Lösern oder PINNs -- eine einzige Vorwärtsberechnung
genügt, um für beliebige neue Eingabefunktionen eine Ausgabe zu produzieren, ohne Neukalibrierung
oder Kollokationspunkte.

Der Autor betrachtet zwei Klassen von Operatorlernaufgaben: (i)~parametrische PDEs, bei denen
aus einem variablen Koeffizienten $a$ die zugehörige Lösung $u = L_a^{-1}f$ gelernt wird, und
(ii)~inverse Koeffizientenidentifikation, bei der aus einer verrauschten Lösung $u$ der
zugrundeliegende Koeffizient $a$ rekonstruiert wird.

% ================================================================
\subsection*{2\quad Methodik: Fourier-Typ versus Galerkin-Typ Aufmerksamkeit}

\paragraph{Grundprinzip.}
Der klassische skalierte Skalarprodukt-Attention-Mechanismus lautet:
\begin{equation}
    \mathbf{z} = \mathrm{Attn}(\mathbf{y}) := (\widetilde{Q}\widetilde{K}^\top)V / n,
\end{equation}
wobei $Q = \mathbf{y}W^Q$, $K = \mathbf{y}W^K$, $V = \mathbf{y}W^V \in \mathbb{R}^{n \times d}$
die Projektionstransformationen sind. Das Paper schlägt vor, die Softmax-Normalisierung vollständig
zu entfernen und stattdessen eine mesh-gewichtete Normalisierung einzusetzen, die eine
Integrations-basierte Interpretation erlaubt.

\paragraph{Fourier-Typ Aufmerksamkeit (FT, quadratische Komplexität).}
\begin{equation}
    \mathbf{z} = \mathrm{Attn}_{\mathrm{f}}(\mathbf{y}) := (\widetilde{Q}\widetilde{K}^\top) V / n.
\end{equation}
Hier werden $Q$ und $K$ vor dem Skalarprodukt mit einer Schicht-Normalisierung versehen.
Die resultierende Ausgabe approximiert ein Fredholm-Integral zweiter Art mit einem
nicht-symmetrischen lernbaren Kernoperator $\kappa(x, \xi) = \zeta_q(x)\phi_k(\xi)$.
Die Komplexität pro Schicht beträgt $\mathcal{O}(n^2 d)$.

\paragraph{Galerkin-Typ Aufmerksamkeit (GT, lineare Komplexität).}
Durch Ausnutzen der Assoziativität der Matrizenmultiplikation ($K^\top V$ zuerst berechnen)
ergibt sich:
\begin{equation}
    \mathbf{z} = \mathrm{Attn}_{\mathrm{g}}(\mathbf{y}) := Q(\widetilde{K}^\top \widetilde{V}) / n.
\end{equation}
Die Komplexität reduziert sich auf $\mathcal{O}(nd^2)$, d.h.\ linear in der Sequenzlänge $n$.
Theoretisch entspricht diese Operation einer lernbaren Petrov-Galerkin-Projektion:
Die Ausgabe $z_j(x)$ approximiert $f$ durch eine Linearkombination von Basisfunktionen
aus dem Spaltenraum von $Q$, wobei die Koeffizienten aus dem inneren Produkt von $K$
und $V$ bestimmt werden (Theorem 4.3, Céa-type Lemma). Diese Approximationskapazität
ist \emph{unabhängig von der Sequenzlänge} $n$ -- ein zentrales theoretisches Ergebnis.

\paragraph{Neue Schicht-Normalisierung.}
Um die Trainingsinstabilität ohne Softmax zu beheben, wird eine diagonale
Galerkin-Projektions-Schicht-Normalisierung eingeführt, bei der $K$ und $V$ (statt $Q$ und $K$)
vor dem Skalarprodukt normalisiert werden. Diese Normalisierung entspricht einer
billigen diagonalen Alternative zu den Normalierungen, die aus dem Petrov-Galerkin-Beweis
folgen, und erlaubt es einem lernbaren Skalierungsfaktor, durch die Encoder-Schichten zu
propagieren.

% ================================================================
\subsection*{3\quad Experimente und Ergebnisse}

Es werden drei Benchmarks untersucht:
\begin{enumerate}
    \item \textbf{Viskose Burgers-Gleichung} (1D, periodisch): Alle aufmerksamkeitsbasierten
    Operatorlerner erzielen vergleichbare Leistung wie FNO1d. Die neue Galerkin-Schicht-Normalisierung
    verbessert das Ergebnis deutlich, wenn Daten unnormalisiert vorliegen.

    \item \textbf{Darcy-Strömung} (2D, elliptisch): Die aufmerksamkeitsbasierten Modelle
    erzielen \textbf{30--50\,\% niedrigeren Evaluierungsfehler} als das Baseline-FNO2d bei
    gleicher Trainingskonfiguration. Der GT Ln on $K,V$ erreicht den besten Wert ($8.39 \times 10^{-3}$
    gegenüber $10.9 \times 10^{-3}$ für FNO2d).

    \item \textbf{Inverse Koeffizientenidentifikation} (2D, Darcy): Für dieses schlecht
    gestellte inverse Problem schlägt das aufmerksamkeitsbasierte Modell das FNO2d deutlich,
    da FNO2d hochfrequente Merkmale durch seine Modus-Trunkierung nicht rekonstruieren kann.
\end{enumerate}

\begin{center}
\begin{tabular}{lrr}
\toprule
\textbf{Modell} & \textbf{CUDA Mem (GB)} & \textbf{GFLOP} \\
\midrule
ST (Softmax) & 31.06 & 1393 \\
FT (Fourier-Typ) & 22.92 & 1138 \\
GT (Galerkin-Typ) & \textbf{1.93} & \textbf{275} \\
\bottomrule
\end{tabular}
\end{center}

Der GT ist damit um einen Faktor $\sim 16\times$ speichereffizienter und $\sim 5\times$
recheneffizienter als der Standard-Softmax-Transformer bei $n = 8192$.

% ================================================================
\subsection*{4\quad Bezug zur Masterarbeit}

\paragraph{4.1\quad Attention als globale Alternative zu FNO-Spektralschichten.}
Der FNO operiert im Spektralbereich und erfasst globale Abhängigkeiten durch Fourier-Koeffizienten.
Die Galerkin-Typ Aufmerksamkeit bietet einen konzeptuell verwandten Zugang: Sie approximiert
das Lösungsfeld durch eine Linearkombination von lernbaren Basisfunktionen und erfasst
globale Korrelationen in einem einzigen Berechnungsschritt. Beide Mechanismen sind für
elliptische Operatoren wie die Stokes-Gleichungen strukturell geeignet, da dort lokale Störungen
das gesamte Rechengebiet beeinflussen.

\paragraph{4.2\quad Relevanz für das Hochfrequenz-Problem an Kristallgrenzen.}
Ein bekanntes Defizit des Standard-FNO ist die Trunkierung hoher Fourier-Moden, die zu
ungenauer Vorhersage an Kristallgrenzen führen kann. Die Galerkin-Typ Aufmerksamkeit
arbeitet dagegen im Ortsraum und hat prinzipiell Zugang zu hochfrequenten räumlichen
Merkmalen -- sofern die Basisfunktionen ausreichend expressiv sind. Dies macht sie zu einem
interessanten Vergleichsmodell für zukünftige Arbeiten, die das Grenzschichtproblem
gezielt adressieren.

\paragraph{4.3\quad Stationäre elliptische PDE: Darcy-Strömung als engste Analogie.}
Die Darcy-Strömung -- $-\nabla \cdot (a \nabla u) = f$, eine stationäre elliptische PDE
zweiter Ordnung -- ist strukturell eng mit der in der Masterarbeit verwendeten
Stokes-Strömung verwandt: Beide sind lineare elliptische Operatoren, bei denen lokale
Inhomogenitäten globale Lösungsstrukturen erzeugen. Die Ergebnisse zeigen, dass
aufmerksamkeitsbasierte Operatorlerner für elliptische Probleme gegenüber FNO
klare Vorteile besitzen -- eine Beobachtung, die die Entwicklung derartiger Modelle
für Stokes-Strömungssurrogate motiviert.

\paragraph{4.4\quad Abgrenzung.}
Das Galerkin-Transformer-Modell operiert auf regulären Gittern und setzt eine einzige
skalare Eingabefunktion (z.B.\ Koeffizientenfeld $a$) voraus. In der Masterarbeit hingegen
werden fünf Eingabekanäle (Kristallmaske, Signed-Distance-Feld, Abstandsfeld, $x$- und
$z$-Koordinaten) auf einem $256 \times 256$-Gitter verwendet, und die geometrische Topologie
variiert zwischen den Proben (unterschiedliche Kristallzahlen). Die direkte Anwendung des
Galerkin-Transformers auf dieses Problem würde Anpassungen an der Eingabekodierung
erfordern. Zudem sind im Galerkin-Transformer-Paper keine Generalisierungsexperimente
auf variierende geometrische Topologien (analog zu variierende Kristallzahlen) enthalten.

\end{document}
